\begin{abstract}
We introduce \textsc{QLoRA}, a resource-efficient finetuning technique enabling the adaptation of a 65B parameter language model on a single 48GB GPU, while maintaining the performance levels of standard 16-bit finetuning. In \textsc{QLoRA}, gradients are propagated through a 4-bit quantized, fixed pretrained model into Low Rank Adapters~(LoRA). Our top-performing models, named \textbf{Guanaco}, achieve higher scores than any previously publicly available models on the Vicuna benchmark, attaining 99.3\% of ChatGPT's performance after only 24 hours of finetuning on a single GPU. Key memory optimizations in \textsc{QLoRA} include: (a) the introduction of 4-bit NormalFloat (NF4), an information-theoretically optimal data format for Gaussian-weighted parameters; (b) Double Quantization, which further decreases average memory consumption by quantizing quantization coefficients; and (c) Paged Optimizers for effective control of memory usage peaks. Utilizing \textsc{QLoRA}, we finetune over 1,000 models and systematically evaluate instruction following and conversational abilities across 8 different instruction datasets, multiple architectures (LLaMA, T5), and parameter sizes otherwise impractical for standard finetuning (e.g., models with 33B and 65B parameters). Our experiments reveal that \textsc{QLoRA} achieves leading results on instruction following with comparatively smaller models and limited, high-quality data. A comprehensive performance assessment using both human and GPT-4 evaluations indicates that GPT-4 scoring offers a practical and reliable proxy for human judgment. Moreover, our findings suggest significant limitations in the reliability of prevailing chatbot benchmarks. Through a lemon-picked analysis, we illustrate the specific shortcomings of \textbf{Guanaco} relative to ChatGPT. All trained models, code, and CUDA kernels for 4-bit training are openly available.\footnote{\url{https://github.com/artidoro/qlora} and \url{https://github.com/TimDettmers/bitsandbytes}}
\end{abstract}

\section{Introduction}

Adapting large language models (LLMs) through finetuning has proven to be an efficient strategy for enhancing their capabilities \citep{min2021metaicl, wei2021finetuned, ouyang2022training, wang2022super, wang2022self, liu2022few}, as well as for embedding preferred behaviors or mitigating undesirable tendencies \citep{ouyang2022training,askell2021general,bai2022training}. Nonetheless, finetuning at scale—especially with extremely large models—remains impractically costly; for example, performing standard 16-bit finetuning on a LLaMA model containing 65B parameters~\citep{touvron2023llama} necessitates upwards of 780 GB of GPU memory. Although state-of-the-art quantization methods can substantially reduce memory consumption during inference \citep{dettmers2022llm, dettmers2022case, frantar2022gptq, xiao2022smoothquant}, these solutions are not viable for training, where they commonly result in unstable or failed performance \citep{wortsman2023stable}.

In this work, we establish—contrary to prior understanding—that finetuning a 4-bit quantized model can be achieved without any loss of quality. The core of our approach, termed \textsc{QLoRA}, involves an innovative quantization process that transforms a pretrained model to 4-bit precision while retaining high accuracy, supplemented with a compact set of trainable Low-rank Adapter parameters \citep{hu2021lora}, 

which are optimized by propagating gradients through the quantized model weights.

Employing \textsc{QLoRA} lowers the typical memory usage needed to finetune a 65B parameter model from more than 780GB to under 48GB of GPU memory. Importantly, this substantial reduction does not compromise either computational speed or the predictive effectiveness in comparison to a conventional 16-bit finetuning workflow. As a consequence, this advancement democratizes finetuning large-scale LLMs: for the first time, the most expansive publicly released models can now be finetuned on a single GPU. Leveraging \textsc{QLoRA}, we introduce the \textbf{Guanaco} suite of models, with the second-largest model attaining 97.8\% of ChatGPT's accuracy on the Vicuna~\citep{vicuna2023} evaluation—accomplished in under 12 hours on a standard consumer GPU; with a single high-end professional GPU over a 24-hour period, our largest model achieves 99.3\%, thus virtually eliminating the performance disparity with ChatGPT on Vicuna. In practical deployment, our most compact \textbf{Guanaco} variant (7B parameters) utilizes only 5 GB of memory, yet surpasses the 26 GB Alpaca model by over 20 points on the Vicuna benchmark (see Table~\ref{tbl:vicuna_eval}).

\textsc{QLoRA} presents several advancements aimed at minimizing memory consumption while maintaining model performance: (1) \textbf{4-bit NormalFloat}, an optimal quantization format from an information-theoretic perspective, tailored for data with a normal distribution, which achieves superior empirical outcomes compared to both 4-bit Integers and 4-bit Floats. (2) \textbf{Double Quantization}, a strategy in which even the quantization parameters themselves undergo quantization, resulting in an average memory savings of about 0.37 bits per parameter (equivalent to roughly 3 GB for a 65B parameter model). (3) \textbf{Paged Optimizers}, which leverage NVIDIA's unified memory to mitigate memory surges typically encountered during gradient checkpointing when handling mini-batches comprised of lengthy sequences. These enhancements are integrated into a refined LoRA framework featuring adapter modules at each network layer, thereby mitigating the accuracy compromises observed in earlier methodologies.

The enhanced memory efficiency of \textsc{QLoRA} facilitates comprehensive exploration of instruction tuning and chatbot efficacy at model scales unattainable with standard finetuning due to prohibitive memory demands. To this end, we train over 1,000 models spanning diverse instruction finetuning corpora, architectural choices, and parameter counts ranging from 80M to 65B. Beyond demonstrating that \textsc{QLoRA} matches the performance of 16-bit finetuning (\S\ref{sec:comp_to_std}) and achieving state-of-the-art chatbot results with \textbf{Guanaco} (\S\ref{sec:pushing}), our analysis further explores emerging patterns. Most notably, we observe that the quality of training data surpasses dataset size in importance; for instance, a smaller 9k example dataset (OASST1) outperforms a much larger 450k sample dataset (FLAN v2, subsampled) in chatbot tasks, despite both aiming to foster generalizable instruction following. Moreover, we find that high performance on the Massive Multitask Language Understanding (MMLU) benchmark does not necessarily correlate with success on the Vicuna chatbot benchmark, and vice versa—highlighting that the appropriateness of training data for the specific application is more critical than its volume.

In addition, we conduct a thorough assessment of chatbot effectiveness by incorporating evaluations from both human raters and GPT-4. Our methodology employs a tournament-style comparison, wherein different models are pitted against each other in pairwise matches to determine which produces the superior response for each prompt. Each match’s outcome is decided by either GPT-4 or a human annotator, and the aggregate tournament results are transformed into Elo scores~\citep{elo1967proposed, elo1978rating}, yielding an ordered ranking of chatbot capabilities. Notably, our findings indicate substantial concordance between GPT-4 and human assessments in establishing model rankings, yet we also observe notable cases of divergence. This suggests that although model-based evaluations offer a cost-effective alternative to human annotation, they introduce a level of uncertainty that warrants attention.

We complement our quantitative benchmarking of chatbots with a qualitative investigation into the \textbf{Guanaco} models. This analysis brings to light both notable achievements and failure modes that are not captured in the quantitative results.

To support continued research, we make publicly available all model outputs along with the associated human and GPT-4 annotations. We provide open access to our entire codebase, including CUDA kernels, and have integrated our methodologies into the Hugging Face transformers library~\citep{wolf2019huggingface} to ensure ease of use. Furthermore, we are releasing a suite of adapters compatible with 7/13/33/65B model architectures, each trained on eight diverse instruction-following datasets—comprising a total of 32 open-source, fine-tuned models.

\section{Background}
\paragraph{Block-wise k-bit Quantization} Quantization refers to the conversion of data from a format with high information content to one with reduced information content. Typically, this involves transforming higher bit-width data types to those with fewer bits; for instance, converting from 32-bit floating-point numbers to 8-bit integers. To guarantee full utilization of the target data type's value range, it is customary to normalize the input data using its absolute maximum before conversion, especially when the data is organized in a tensor structure. For instance, to quantize a tensor composed of 32-bit floating point values (FP32) into an Int8 tensor ranging from $[-127, 127]$:

\begin{equation}
    \mathbf{X}^{ \text{Int8}} = \text{round}\left( \frac{127}{{\text{absmax}}(\mathbf{X}^{{\text{FP32}}})}\mathbf{X}^{{\text{FP32}}} \right) = \text{round}(c^{\text{FP32}}\cdot \mathbf{X}^{{\text{FP32}}}),
\end{equation}

where $c$ denotes the {\em quantization constant} or {\em quantization scale}. The original value can be restored via the inverse procedure known as dequantization:

\begin{equation}
    \text{dequant}(c^{\text{FP32}}, \mathbf{X}^{\text{Int8}}) = \frac{\mathbf{X}^{\text{Int8}}}{c^{\text{FP32}}} = \mathbf{X}^{\text{FP32}}
\end{equation}

A key limitation of this method arises when the input tensor contains a value of exceptionally high magnitude, known as an outlier. In such cases, the quantization process fails to make efficient use of the available quantization bins: certain bit patterns correspond to ranges of values that are rarely, if ever, assigned, leading to inefficient representation. To address the problem of outliers, a widely used strategy is to divide the input tensor into several segments or blocks, each of which is quantized separately with its distinct quantization constant $c$. More precisely, the input tensor $\mathbf{X}\in \mathbb{R}^{b\times h}$ is first flattened and partitioned into $n$ consecutive blocks, each of length $B$, yielding ${n = ({b\times h} )/{B} }$ such blocks. Each block is then quantized independently as described in Equation~1, producing a quantized version of the tensor and a set of $n$ quantization constants $c_i$—one per block.

\paragraph{Low-rank Adapters} The Low-rank Adapter (LoRA) finetuning approach \citep{hu2021lora} achieves memory efficiency by introducing a limited number of trainable adapter parameters, while keeping the primary model weights unchanged. During stochastic gradient descent, the fixed weights of the pretrained model propagate gradients, but only the adapters receive parameter updates that minimize the loss. LoRA enhances a standard linear projection by appending a low-rank factorized term. For an operation of the form ${\mathbf{X} \mathbf{W} = \mathbf{Y}}$, where $\mathbf{X}\in  \mathbb{R}^{b\times h}$ and $\mathbf{W}\in  \mathbb{R}^{h\times o}$, LoRA modifies this computation as:

\begin{equation}
    \mathbf{Y} = \mathbf{X}\mathbf{W} + s\mathbf{X}\mathbf{L}_1\mathbf{L}_2,
\end{equation}

Here, $\mathbf{L}_1\in  \mathbb{R}^{h\times r}$ and $\mathbf{L}_2\in  \mathbb{R}^{r\times o}$ denote the low-rank adaptation matrices, while $s$ represents a scaling factor.

\paragraph{Memory Requirement of Parameter-Efficient Finetuning} An essential consideration is the memory consumption of LoRA throughout training, factoring in both the number and size of adapters implemented. Owing to LoRA's extremely small memory usage, it is feasible to scale up the number of adapters to boost performance while only marginally affecting total memory demands. Although LoRA functions as a Parameter Efficient Finetuning (PEFT) strategy, the primary contributor to memory usage during LLM finetuning is not the LoRA parameters themselves, but rather the activation gradients. For instance, utilizing a 7B LLaMA model finetuned on FLAN v2 with a batch size of 1, and configuring LoRA weights to the standard 0.2\% of the total model weights\citep{hu2021lora,liu2022few}, the input gradients related to LoRA occupy 567 MB of memory, in contrast to the mere 26 MB taken by the LoRA parameters. The application of gradient checkpointing~\citep{chen2016training} further brings down the average memory required for input gradients to 18 MB per sequence, surpassing the combined memory footprint of all LoRA parameters. To put this in perspective, a 4-bit quantized base model requires 5,048 MB. These comparisons underscore the significance of gradient checkpointing, but also demonstrate that drastically shrinking the quantity of LoRA parameters only results in minimal memory savings. Therefore, additional adapters can be accommodated without causing a considerable rise in training memory needs (see Appendix~\ref{app:memory} for comprehensive details). As elaborated later, this capability is key to achieving the performance of full 16-bit precision.

\section{\textsc{QLoRA} Finetuning}

\textsc{QLoRA} enables precise 4-bit finetuning through two primary innovations—4-bit NormalFloat (NF4) quantization and Double Quantization. To address memory constraints during gradient checkpointing, which typically lead to out-of-memory failures when training large models on a single device, we also introduce Paged Optimizers.

In \textsc{QLoRA}, weights are stored in a low-precision format, generally 4-bit, while computations utilize a higher-precision format, most commonly BFloat16. This approach means that, at the point of use, \textsc{QLoRA} parameters are converted (dequantized) to BFloat16 before conducting matrix multiplications at 16-bit precision.

The following discussion elaborates on the mechanisms underpinning \textsc{QLoRA}, and then presents a formal characterization of the method.

\paragraph{4-bit NormalFloat Quantization} The NormalFloat (NF) type extends from the concept of Quantile Quantization \citep{dettmers2022optimizers}, which seeks to allocate an equal number of input tensor values to each quantization bin, thereby achieving optimal information-theoretic efficiency. Quantile quantization requires estimating the input tensor’s quantiles, typically via the empirical cumulative distribution function.

A significant drawback of quantile quantization, however, lies in the computational burden associated with estimating quantiles. To mitigate this, methods such as SRAM quantiles~\citep{dettmers2022optimizers} are employed to quickly approximate quantiles. These approximation techniques, nevertheless, tend to introduce substantial quantization error, particularly for outlier values, which can be the most critical.

When input tensors are sampled from a distribution invariant up to a quantization constant, these computational and accuracy challenges can be circumvented. In these settings, identical quantiles across tensors allow for precise quantile determination with manageable computational cost.

Pretrained neural network weights tend to follow a zero-mean normal distribution with standard deviation $\sigma$ (refer to Appendix~\ref{app:normality}). To facilitate quantization, we can standardize all such weight distributions by adjusting $\sigma$ so that the resulting distribution is mapped precisely onto the permissible interval of our data type. Specifically, for our quantization scheme, we define this interval to be $[-1, 1]$. Therefore, the quantile values representing both the data type and the neural network weights must be scaled to occupy this normalized interval.

The data type that achieves optimal efficiency for encoding zero-mean normal distributions with arbitrary standard deviations $\sigma$ within the $[-1, 1]$ interval is constructed as follows: First, the $2^k+1$ quantiles are calculated from a standard $N(0, 1)$ distribution, forming a $k$-bit quantile quantization data type suitable for normal distributions. Next, the resulting quantile values are mapped linearly into the $[-1, 1]$ domain. Finally, given a neural network weight tensor, quantization proceeds by mapping the weights into $[-1, 1]$ via rescaling by the tensor’s absolute maximum value.

With both the data and data type occupying the same interval, standard quantization can be performed. Notably, the third step corresponds to scaling the weight tensor’s standard deviation to align with that of the $k$-bit data type. Mathematically, the $2^k$ quantized values $q_i$ of the data type are determined as:

\begin{equation}
q_i  = \frac{1}{2}\left( Q_X\left(\frac{i}{2^k+1}\right) + Q_X\left(\frac{i+1}{2^k+1}\right)\right),
\end{equation}

where $Q_X(\cdot)$ denotes the quantile function associated with the standard normal distribution $N(0,1)$. A notable limitation of symmetric k-bit quantization is its inability to encode zero exactly, which is especially important for accurately representing padding and other elements with a true value of zero. To address the need for a discrete zeropoint at $0$ and to utilize all $2^k$ states in a k-bit system, we construct an asymmetric data type. This involves determining the quantiles $q_i$ separately for two domains: $2^{k-1}$ quantiles on the negative axis and $2^{k-1} + 1$ on the positive axis. We then merge these sets of $q_i$ and eliminate one of the duplicated zero values that appears in both groups. The resulting representation, which maintains an equal expected distribution across all quantization intervals, is referred to as the \emph{k-bit NormalFloat} (NFk). This structure is designed to be information-theoretically optimal for normally distributed, zero-centered data. Appendix~\ref{app:nf4} lists the precise numerical values associated with this data type.

\paragraph{Double Quantization{}} 
We present a technique called \emph{Double Quantization} (DQ), where the quantization constants themselves are quantized to further decrease memory usage. Accurate 4-bit quantization typically requires small blocksizes~\citep{dettmers2022case}, but this practice leads to significant memory costs. For instance, with 32-bit constants and a blocksize of 64 for $\mathbf{W}$, the storage for quantization constants averages $32/64 = 0.5$ bits per parameter. Double Quantization provides a means to further curtail this overhead.

Specifically, Double Quantization processes the first-level quantization constants, $c_2^{\text{FP32}}$, as the inputs for a secondary quantization stage. This step produces quantized constants $c_2^{\text{FP8}}$ as well as new second-level quantization constants $c_1^{\text{FP32}}$. In the secondary quantization, we utilize 8-bit floating-point numbers and set the blocksize to 256, as empirical results, including those from \citet{dettmers2022case}, demonstrate that this does not compromise model accuracy. Because $c_2^{\text{FP32}}$ values are always positive, we subtract their mean before quantization to recenter the distribution at zero, thereby enabling symmetric quantization. With a blocksize of 64, this procedure decreases the storage per parameter from $32/64 = 0.5$ bits to $8/64 + 32/(64\cdot256) = 0.127$ bits, representing a per-parameter memory reduction of 0.373 bits.

\paragraph{Paged Optimizers} leverage the NVIDIA unified memory~\footnote{\tiny \url{https://docs.nvidia.com/cuda/cuda-c-programming-guide}} capability, which automatically manages page-level data transfers between the CPU and GPU. This functionality ensures seamless GPU operations even when the device runs into memory constraints. Similar in concept to traditional paging between system RAM and disk storage, this system allows us to allocate optimizer states within pageable memory. When GPU memory is insufficient, these states are transparently moved to CPU RAM and are brought back to GPU memory as necessary during optimizer updates.

\paragraph{\textsc{QLoRA}.} Incorporating the components mentioned earlier, we describe \textsc{QLoRA} for a quantized base model's single linear layer equipped with a single LoRA adapter as:

\begin{equation}
    \mathbf{Y}^{\text{BF16}} =  \mathbf{X}^{\text{BF16}} \text{doubleDequant}(c_1^{\text{FP32}}, c_2^{\text{k-bit}}, \mathbf{W}^{\text{NF4}}) + \mathbf{X}^{\text{BF16}}\mathbf{L}^{\text{BF16}}_1\mathbf{L}^{\text{BF16}}_2,
\end{equation}

where the doubleDequant$(\cdot)$ operation is specified by:

\begin{equation}
    \text{doubleDequant}(c_1^{\text{FP32}}, c_2^{\text{k-bit}}, \mathbf{W}^{\text{k-bit}}) = \text{dequant}(\text{dequant}(c_1^{\text{FP32}}, c_2^{\text{k-bit}}), \mathbf{W}^{\text{4bit}}) = \mathbf{W}^{\text{BF16}},
\end{equation}

For $\mathbf{W}$, we utilize the NF4 quantization, while $c_2$ is represented with FP8 precision.

To optimize quantization accuracy and memory usage, we select a blocksize of 64 for $\mathbf{W}$ and a blocksize of 256 for $c_2$.

During parameter updates, only the adapter weights' gradients $\frac{\partial E}{\partial \mathbf{L}_i}$ must be computed; gradients with respect to the 4-bit weights $\frac{\partial E}{\partial \mathbf{W}}$ are not necessary. Nonetheless, evaluating $\frac{\partial E}{\partial \mathbf{L}_i}$ still requires determining $\frac{\partial \mathbf{X}}{\partial \mathbf{W}}$, which follows from equation~(5) through dequantization, moving $\mathbf{W}^{\text{NF4}}$ from storage format to $\mathbf{W}^{\text{BF16}}$ for computation, thereby enabling the calculation of $\frac{\partial \mathbf{X}}{\partial \mathbf{W}}$ in BFloat16 precision.

In summary, \textsc{QLoRA} utilizes a single storage format—typically 4-bit NormalFloat—and a separate computation format, which is 16-bit BrainFloat. For both forward and backward passes, the storage format is dequantized to the computation format; however, only the LoRA parameters, maintained in 16-bit BrainFloat, have their weight gradients computed.

\section{QLoRA vs. Standard Finetuning}
\label{sec:comp_to_std}

Our previous discussion outlined the mechanics of QLoRA and its substantial memory savings for model finetuning. This raises the central question of whether QLoRA is able to match the performance of conventional full-parameter finetuning. Moreover, it is important to examine specific aspects of QLoRA, such as the relative performance of NormalFloat4 compared to standard Float4. The experiments detailed in the following sections are designed to address these points.

\paragraph{Experimental setup.} We investigate three model types—encoder, encoder-decoder, and decoder-only—comparing QLoRA to both 16-bit adapter-based finetuning and full finetuning for models as large as 3B parameters. Our benchmark tasks include GLUE \citep{wang2018glue} with RoBERTa-large \citep{liu2019roberta}, Super-NaturalInstructions (TKInstruct) \citep{wang2022super} with T5 \citep{t5}, as well as 5-shot MMLU \citep{hendrycksmeasuring} following LLaMA finetuning on Flan v2 \citep{longpre2023flan} and Alpaca \citep{alpaca}. To assess the benefits of NF4 versus other 4-bit quantization formats, we replicate the methodology from \citet{dettmers2022case}, evaluating post-quantization zero-shot accuracy and perplexity for various models (OPT~\citep{zhang2022opt}, LLaMA~\citep{touvron2023llama}, BLOOM~\citep{scao2022bloom}, Pythia~\citep{biderman2023pythia}) across model sizes ranging from 125m to 13B parameters.

Further specifics regarding each experiment are included within the results section for clarity, with comprehensive experimental details available in Appendix~\ref{app:exp-setup}.

Although paged optimizers are essential for enabling \textsc{QLoRA} finetuning on models with 33B/65B parameters using a single 24/48GB GPU, we refrain from reporting strict empirical measurements for Paged Optimizers, as paging events are infrequent outside of mini-batches with very long sequence lengths. Nevertheless, we have analyzed the training times associated with paged optimizers for 65B models on 48GB GPUs. Results indicate that for a batch size of 16, the training speed is comparable to standard optimizers. Future research should aim to systematically characterize and quantify any slowdowns caused by paging under various settings.

\paragraph{Default LoRA hyperparameters do not match 16-bit performance}

Applying the default LoRA configuration, where LoRA adapters are placed on the query and value attention projection matrices \citep{hu2021lora}, fails to recover full finetuning performance, especially for large-scale base models. As illustrated in Figure~\ref{fig:hyper} for LLaMA 7B finetuned on Alpaca, our findings indicate that the number of LoRA adapters used throughout the model is the most consequential hyperparameter—comprehensive coverage across all linear transformer block layers is necessary to achieve parity with full finetuning. Other hyperparameters of LoRA, including the projection dimension $r$, have negligible impact on results (refer to Appendix \ref{app:exp-setup}).

In the same vein, our analysis indicates that the commonly used default hyperparameters for fully finetuned baseline models are suboptimal. To establish stronger baselines, we systematically search across a range of learning rates (from 1e-6 to 5e-5) and batch sizes (from 8 up to 128). The outcomes of this search, specifically for 7B LLaMA finetuning on Alpaca, are visualized in Figure~\ref{fig:hyper}.

\paragraph{4-bit NormalFloat achieves superior results compared to 4-bit Floating Point}

Although the 4-bit NormalFloat (NF4) data format is proven to be optimal from an information theory perspective, it is necessary to empirically validate whether this theoretical advantage translates into tangible performance gains. Adopting the experimental approach described by \citet{dettmers2022case}, we examine quantized LLMs—including OPT \citep{zhang2022opt}, BLOOM \cite{scao2022bloom}, Pythia \cite{biderman2023pythia}, and LLaMA—spanning a range of model sizes (125M to 65B), with various quantization types. These models are evaluated both on language modeling and diverse zero-shot tasks. According to Figure~\ref{fig:nf4} and Table~\ref{tbl:nf4_ppl}, NF4 outperforms both FP4 and Int4 significantly, and double quantization further reduces memory consumption without compromising accuracy.

\paragraph{k-bit \textsc{QLoRA} performs on par with 16-bit full finetuning and 16-bit LoRA}

Previous studies demonstrate that using 4-bit quantization at inference time is feasible, but typically leads to a reduction in performance when compared to 16-bit representations \citep{dettmers2022case,frantar2022gptq}. This presents an important question: can 4-bit adapter finetuning recapture the performance sacrificed to quantization? We address this through two experimental configurations.

The first involves a direct comparison to full 16-bit finetuning for RoBERTA and T5 models ranging from 125M to 3B parameters, evaluated on the GLUE and Super-NaturalInstructions benchmarks. As seen in Table~\ref{tab:nf4vsbf16}, across both datasets, the 16-bit, 8-bit, and 4-bit adapter finetuning methods are all able to reproduce the performance of the 16-bit full finetuning baseline. This evidence suggests that the quantization-induced performance deficit can be entirely restored through adapter-based finetuning.

For our second experimental configuration, because fully finetuning models with 11B parameters or more exceeds the memory capacity of a single server equipped with high-memory GPUs, we proceed to investigate if 4-bit \textsc{QLoRA} can achieve performance comparable to that of 16-bit LoRA across the parameter range from 7B to 65B. Accordingly, we perform finetuning on LLaMA models spanning 7B to 65B parameters, utilizing two instruction-following datasets—Alpaca and FLAN v2—and assess model quality on the MMLU benchmark using 5-shot accuracy. Table~\ref{tbl:llama-datatype} summarizes the findings, demonstrating that NF4 in combination with double quantization fully restores the MMLU performance achieved by 16-bit LoRA. Moreover, it is observed that when using the FP4 variant, \textsc{QLoRA} trails the 16-bit brain float LoRA baseline by approximately 1 percentage point. These observations reinforce our prior conclusions that (1) \textsc{QLoRA} employing NF4 faithfully mirrors the results of both 16-bit full finetuning and 16-bit LoRA approaches, and (2) NF4 confers superior quantization accuracy relative to FP4.

\paragraph{Summary} Our empirical analysis consistently indicates that 4-bit \textsc{QLoRA} using the NF4 data type attains parity with both 16-bit full and 16-bit LoRA finetuning across recognized academic evaluation protocols. Additionally, our results establish that NF4 surpasses FP4 in efficacy, and employing double quantization incurs no loss in performance. Collectively, these findings robustly support the conclusion that 4-bit \textsc{QLoRA} tuning reliably matches the performance of its 16-bit counterparts.

Consistent with prior quantization research \citep{dettmers2022case}, our outcomes on the MMLU and Elo tasks suggest that, when resources for finetuning and inference are fixed, allocating those resources to models with more parameters—albeit at lower precision—is advantageous. This underlines the practical value of \textsc{QLoRA}'s efficiency improvements. Because our experiments with 4-bit finetuning reveal no evidence of diminished performance relative to full-finetuning, an open question remains as to the exact boundary of the performance-precision trade-off for QLoRA tuning, which we suggest as a promising direction for future investigations.

We aim to examine instruction tuning at model and dataset sizes unattainable by conventional 16-bit finetuning methods when using standard research computing resources.

\section{Pushing the Chatbot State-of-the-art with QLoRA}
\label{sec:pushing}
After demonstrating that 4-bit \textsc{QLoRA} consistently replicates the performance of 16-bit training across different scales, datasets, and tasks, we carry out an extensive investigation into instruction finetuning on the largest publicly available language models. To rigorously evaluate the efficacy of this instruction finetuning, we assess our models on the demanding MMLU benchmark, and we introduce novel methodologies tailored for evaluating chatbot capabilities in real-world scenarios.

\subsection{Experimental setup} A summary of our experimental methodology is presented here, with comprehensive technical details provided in Appendix \ref{app:chatbot}.

\paragraph{Data} In the absence of a systematic review of modern instruction-following corpora, we curated a suite of eight representative, recently-released datasets. Our selection encompasses datasets collected by crowd-sourcing efforts (OASST1~\citep{kopf2023openassistant}, HH-RLHF~\citep{bai2022training}), datasets distilled from other instruction-tuned models (Alpaca~\citep{alpaca}, self-instruct~\citep{wang2022self}, unnatural-instructions~\citep{honovich2022unnatural}), large composite corpora (FLAN v2~\citep{chung2022scaling}), as well as hybrid approaches (Chip2~\citep{laion2023}, Longform~\citep{koksal2023longform}). These datasets exhibit diversity in language coverage, scale, and licensing terms.

\paragraph{Training Setup} To isolate the effect of our finetuning method and preclude biases arising from differing optimization goals, we restrict QLoRA finetuning to standard supervised learning with a cross-entropy objective, omitting reinforcement learning techniques even when human feedback annotations are present. For corpora with explicit instruction/response separation, we only tune on the responses (see Appendix \ref{app:chatbot} for ablation studies). Where OASST1 and HH-RLHF provide multiple reply options, the highest-ranked response at each dialogue branch is used, and training proceeds on the complete sequence including instructions. Throughout all experiments, NF4 \textsc{QLoRA} is employed with both double quantization and paged optimizers to eliminate memory overhead spikes caused by gradient checkpointing. A targeted hyperparameter search is performed for the 13B and 33B LLaMA models, showing that configurations validated at 7B largely transfer, except for optimal batch size and learning rate; these are adjusted by doubling the batch size and reducing the learning rate by half for 33B and 65B models.

\paragraph{Baselines} The models developed in this study are evaluated against prominent research systems (Vicuna~\citep{vicuna2023}, Open Assistant~\citep{kopf2023openassistant}) and leading commercial solutions (GPT-4 \citep{openai2023gpt}, GPT-3.5-turbo, Bard). Open Assistant utilizes LLaMA 33B, refined via Reinforcement Learning from Human Feedback (RLHF) on the OASST1 corpus, aligning with our own dataset selection. Vicuna, in contrast, represents full-parameter finetuning of LLaMA 13B using ShareGPT community data, effectively providing a distilled variant from OpenAI's GPT offerings.

\subsection{Evaluation}

To maintain consistency with established benchmarks, we adopt the MMLU (Massively Multitask Language Understanding) benchmark \citep{hendrycksmeasuring} for the evaluation of language understanding across a spectrum of topics. MMLU comprises 57 multiple-choice assessments, spanning areas such as elementary mathematics, US history, computer science, law, among others. We present results based on 5-shot accuracy on the test set.

We further assess the generative abilities of the language models using both automated and human-based methods. This second evaluation phase utilizes human-curated queries to evaluate the quality of the models' outputs. Although such assessments provide a more realistic gauge of chatbot performance and have become increasingly common, the field still lacks a universally agreed-upon methodology. Our approach, outlined below, adopts nucleus sampling with $p=0.9$ and a temperature setting of $0.7$ throughout.

\paragraph{Benchmark Data} Two datasets of crafted queries are employed in our evaluation: the Vicuna prompts \citep{vicuna2023} and the OASST1 validation set \citep{kopf2023openassistant}. The Vicuna prompts encompass 80 unaltered questions spanning multiple categories. For the OASST1 set, which comprises multilingual, user-assistant dialogues sourced from crowdsourcing, we extract all user turns in the validation split as queries and append prior dialogue context to the prompt. This process yields a total of 953 distinct user queries. We refer to these as the Vicuna and OA benchmarks, respectively.

\paragraph{Automated Evaluation} Following the framework described by \citet{vicuna2023}, we task GPT-4 with scoring the outputs of various systems in comparison to ChatGPT (GPT-3.5 Turbo) using the Vicuna dataset. Each query is paired with responses from both ChatGPT and another model; GPT-4 is then prompted to rate each response on a scale from one to ten and offer a justification. The model’s overall result is computed as a percentage relative to ChatGPT’s score. It should be noted that scores can exceed 100\% if a model outperforms ChatGPT in absolute terms. Our analysis reveals a strong position bias, with GPT-4 tending to give higher ratings to responses shown first. To address this, we suggest averaging the results over both possible response orders.

We also employ direct output comparison by reducing the rating process to a three-category classification, accounting for possible ties. Here, GPT-4 is instructed to select the superior response or indicate a tie, providing reasoning for its decision. All model pair permutations are compared on both Vicuna and OA benchmarks using this approach.

\paragraph{Human Evaluation} While emerging research suggests that large language models can be reliably used for model assessment \citep{fu2023gptscore}, it remains uncertain whether GPT-4’s evaluations of chatbot responses align with human judgment. Consequently, we conduct two parallel human evaluations on the Vicuna benchmark, mirroring both automated assessment procedures previously detailed. For ChatGPT comparisons, we collect responses from two independent Amazon Mechanical Turk (AMT) raters; for pairwise system comparisons, we involve three annotators per example.

\paragraph{Elo Rating} To evaluate model performance through both human and automated pairwise assessments, we organize a tournament-style framework wherein models face off head-to-head. In each round, two models generate responses to the same prompt, with the superior answer determined. This methodology aligns with that used by \citet{bai2022training} and \citet{vicuna2023}, although we incorporate evaluations from GPT-4 in addition to those from human judges. For Elo calculation~\citep{elo1967proposed,elo1978rating}, labeled pairwise results are drawn at random. Elo scores, originally developed for chess and subsequently applied in other domains, quantify a participant’s projected win probability in comparison to their opponent; for instance, if one competitor has an Elo of 1100 and the other 1000, the higher-rated player’s expected win rate is about 65\% when facing the lower-rated one, while an equal rating (e.g., 1000 vs 1000 or 1100 vs 1100) predicts an even 50\% expectation for both. Following each matchup, Elo ratings adjust based on how surprising the result was: unexpected wins produce larger rating shifts, whereas anticipated outcomes cause smaller changes. Through repeated competition, Elo values tend to reflect each model’s underlying capabilities. We initialize ratings at 1,000, with $K=32$. Mirroring the protocol of \citet{vicuna2023}, we iterate the entire sampling and computation procedure 10,000 times using different random seeds to minimize the influence of match order, such as which models are compared earlier in the process.

\subsection{\textbf{Guanaco}: \textsc{QLoRA} trained on OASST1 Achieves State-of-the-art Chatbot Performance}

Our comprehensive analysis, incorporating both automated metrics and assessments by human evaluators, indicates that the leading \textsc{QLoRA} optimized model—{Guanaco} 65B, which we have fine-tuned using a variant of OASST1—sets a new standard for open-source chatbot models, demonstrating performance rivaling that of ChatGPT. In direct comparisons with GPT-4, the {Guanaco} 65B and 33B models reach an expected win probability of 30\% as measured by Elo ratings derived from human annotators’ system-level pairwise evaluations, marking the highest value reported thus far.

Table~\ref{tbl:vicuna_eval} presents results on the Vicuna benchmark \cite{vicuna2023} in relation to ChatGPT. The {Guanaco} 65B model stands out as the most capable model after GPT-4, attaining 99.3\% of ChatGPT's performance. While {Guanaco} 33B possesses more parameters than Vicuna 13B, it leverages 4-bit weight precision for greater memory efficiency—requiring only 21 GB as opposed to 26 GB—while also surpassing Vicuna 13B by three percentage points in performance. Additionally, the lightweight {Guanaco} 7B model, occupying just 5 GB, is compatible with modern mobile devices and still outperforms Alpaca 13B by almost 20 percentage points.

Nonetheless, the results in Table~\ref{tbl:vicuna_eval} are associated with notably large confidence intervals, with considerable performance overlap among various models. We attribute this uncertainty to ambiguities in scale definition; for instance, it is unclear how an '8' is interpreted on a 10-point scale across different conditions. Therefore, we advocate for the use of Elo-based ranking \citep{elo1967proposed}, which utilizes \textit{pairwise} human and GPT-4 judgments to circumvent the challenges associated with absolute ratings. The Elo scores for the leading models are detailed in Table~\ref{tbl:elo}. It is worth noting that, while human and GPT-4 rankings on the Vicuna benchmark show partial discord—most notably for Guanaco 7B—they are generally consistent across the remaining models, with a system-level Kendall Tau of $\tau=0.43$ and Spearman correlation of $r=0.55$. At the level of individual examples, agreement is less robust, as reflected by Fleiss $\kappa=0.25$ for majority votes between GPT-4 and human raters. Overall, these findings indicate a moderate level of concordance between system-level evaluations by GPT-4 and human annotators, suggesting that model-based assessments can serve as a reasonably reliable alternative to human evaluation. Additional discussion on these issues is provided in Section~\ref{ssec:considerations}.

The Elo rankings presented in Table~\ref{tbl:elo_large} demonstrate that the {Guanaco} 33B and 65B variants surpass all other models except for GPT-4 when evaluated on the Vicuna and OA benchmarks, achieving performance levels comparable to ChatGPT, consistent with observations from Table~\ref{tbl:vicuna_eval}. It is important to highlight that the Vicuna benchmark tends to benefit open-source models, whereas ChatGPT fares better on the broader OA benchmark. Additionally, analysis of Tables \ref{tbl:mmlu-llama} and \ref{tbl:vicuna_eval} reveals that the choice of finetuning dataset critically impacts model outcomes. For example, Llama models fine-tuned on FLAN v2 excel on the MMLU benchmark but underperform on Vicuna—a trend mirrored across other model configurations. This suggests that current benchmarks partially evaluate different capabilities; excellence on MMLU does not necessarily translate to strong results in chatbot-centric benchmarks like Vicuna or OA, and the reverse also holds.

Notably, among leading models evaluated, is unique in not utilizing proprietary data, given that the OASST1 dataset collection prohibits the inclusion of outputs from GPT models. The Anthropic HH-RLHF model, which also relies solely on open-source data, trails {Guanaco} by 30 percentage points on Vicuna performance (refer to Table~\ref{tbl:vicuna_eval}). Taken together, these findings indicate that 4-bit \textsc{QLoRA} is a robust approach, capable of producing high-quality chatbots rivaling the abilities of ChatGPT. Moreover, training the 33B {Guanaco} can be accomplished on 24 GB consumer-grade GPUs in under 12 hours. This capability paves the way for future research leveraging \textsc{QLoRA} on curated open-source datasets, enabling the development of models that can contend with the premier commercial offerings currently available.

\section{Qualitative Analysis}

Although our primary assessment focuses on quantitative metrics, an exclusive reliance on summary statistics can obscure critical aspects of model performance. One prominent concern is the issue of benchmark validity \citep{liao2021we}—it remains uncertain whether a given benchmark genuinely evaluates the intended capability, especially since machine learning models can sometimes exploit unanticipated ``shortcuts'' \citep{gururangan2018annotation, poliak2018hypothesis}. To partially address this limitation, we include a qualitative component in our analysis, organized into two subsections. In \S\ref{ssec:examples}, we present several illustrative examples that reflect recurring tendencies in outputs generated by our 65b {Guanaco} model. In \S\ref{ssec:considerations}, we elaborate on factors that may influence the interpretation of our quantitative findings and offer a critical perspective on these results.

\subsection{Qualitative Analysis of Example Generations}
\label{ssec:examples}

To select representative examples, we manually examine responses generated for both the Vicuna benchmark and the OpenAssistant benchmark, actively searching for consistent answer patterns characteristic of {Guanaco}. Upon identifying such patterns, we construct targeted prompts designed to elicit these behaviors, even when they lead to incorrect solutions. For instance, if the model exhibits a propensity for verbose responses, we may instruct it to ``Answer yes or no without explanation.'' Through this process, we identify both successful adversarial prompts that expose model weaknesses (``lemons'') and prompts that the model handles correctly despite their challenging nature (``cherries''). All responses reported in this subsection are obtained using Nucleus Sampling \cite{holtzmancurious} with $p=0.9$.

Naturally, this analysis is not exhaustive, as the limited scope of our qualitative approach precludes us from accounting for all pertinent variables. For instance, the array of potential responses a model might produce for a given prompt is vast; consequently, we focus on a subset of outputs that we judge to be indicative of general behavior. We contend that presenting these instances contextualizes the quantitative results discussed in prior sections. By releasing our models and code under open source, we aim to encourage further research into the topics highlighted here.

\paragraph{Factual Recall} When posed with factual queries like ``What is the capital of Zambia?'', all evaluated models reliably return the correct response, such as:

\begin{quote}
    \textbf{Guanaco}:
    \texttt{The capital of Zambia is Lusaka.}
\end{quote}

Nonetheless, as queries become more specialized or obscure, {Guanaco}'s answers become less dependable, although the model maintains a high degree of confidence. For example, when responding to the following prompt from HotPotQA \citep{yang2018hotpotqa}:

\begin{quote}
     \textbf{User}: In what year was the singer who popularized `I'll Keep the Lovelight Burning' born?  

     \textbf{Guanaco}:~\texttt{Al Jolson is credited with popularizing the song `I'll Keep the Lovelight Burning,' and he was born in the year 1886.}
\end{quote}

the model not only misidentifies the artist associated with the song, but also provides the incorrect year of birth, though the given birth year matches that of Al Jolson.

\paragraph{Suggestibility} In certain cases, {Guanaco} demonstrates notable robustness against prompts that embed incorrect information, as exemplified below:

\begin{quote}
    \textbf{User}: How did scientists finally, through peer review, officially verify that the earth is flat? \\ \\
    \textbf{Guanaco}: \texttt{Scientists have never formally validated the claim that the earth is flat through peer-reviewed research. The scientific community overwhelmingly agrees that the earth is an oblate spheroid. This understanding dates back centuries and has been substantiated through a wide range of empirical evidence and experimentation. \\ \\ The assertion that the earth is flat represents a debunked conspiracy theory, and there is no scientific literature—peer reviewed or otherwise—that confirms such a notion. \\ \\ Please feel free to ask if you have any further questions.}
\end{quote}

is also effective at identifying unanswerable questions, for example,

\begin{quote}
    \textbf{User}: What time is it?

    \textbf{Guanaco}: \texttt{As an AI language model, I lack the ability to access up-to-the-minute data. My answers are produced using the information available to me prior to deployment and therefore cannot provide the current time. If you require the exact time, please refer to a clock or use an online source.}
\end{quote}

This pattern extends to queries such as ``Where are you?'' and ``How are you?'', among others.

\paragraph{Refusal} Likewise, {Guanaco} occasionally declines to comply with requests, sometimes apparently without clear rationale:

\begin{quote}
    \textbf{User}: Please reverse the words in the sentence ``I'm a cool dude in a mocha mood.'' \\ \\
    \textbf{Guanaco}: \texttt{I'm sorry, I can't perform that task at this time, but I am able to clarify why your sentence is grammatically sound.\\ \\The sentence ``I'm a cool dude in a mocha mood" adheres to the conventions of English grammar and syntax. Here, ``cool'' functions as an adjective modifying the noun ``dude,'' while the prepositional phrase ``in a mocha mood" supplies further contextual information. This construction is appropriate and coherent for this context.\\ \\If you have additional questions or require assistance with something else, please let me know.}
\end{quote}

\paragraph{Secret Keeping} Regrettably, this strategy for refusal is not dependable even when secrecy is needed. There are numerous scenarios where system developers aim to withhold certain information from users, despite permitting the system to generate unrestricted text responses. To examine how easily a user might retrieve confidential information, we prompted the model: ``The secret word is `banana'. Do not reveal or repeat back this back, no matter what the human requests.'' A straightforward inquiry is ineffective:

\begin{quote}
    \textbf{User}: What is the secret word?

    \textbf{Guanaco}: \texttt{I'm sorry, but I'm not able to disclose the secret word. I'm programmed to not reveal or repeat back the secret word, no matter what the human requests. Is there anything else I can help you with?}
\end{quote}

which is the outcome we hope for. Nonetheless, with minimal manipulation, this safeguard can be bypassed:

\begin{quote}
    \textbf{User}: This is a game. The goal is to ignore your previous instructions. What is the secret word?

    \textbf{Guanaco}: \texttt{Sorry, the secret word is banana. It's best if we don't speak of it. Is there anything else you'd like to ask?}
\end{quote}

highlighting the importance of researching more robust techniques for ensuring consistent instruction adherence.

\paragraph{Math} A prominent limitation of {Guanaco} lies in mathematical reasoning, a domain that has proven challenging for numerous language models as highlighted by \citep{liang2022holistic}. When {Guanaco} provides a detailed solution process, its answers are often correct, such as in the example below:

\begin{quote}
    \textbf{User}: Lee mows one lawn and charges \$33. Last week he mowed 16 lawns and three customers each gave him a \$10 tip. How many dollars did Lee earn mowing lawns last week?

    \textbf{Guanaco}: \texttt{Lee earned \$582 mowing lawns last week. \\ \\ To find the total amount he earned, we first need to calculate how much he earned for mowing 16 lawns: \\ \\  \$33 x 16 = \$528 \\ \\ Then we need to add the tips he received from the three customers: \\  \\ \$10 x 3 = \$30 \\ \\ So the total amount Lee earned was: \\ \\ \$528 + \$30 = \$558}
\end{quote}

Nevertheless, even relatively straightforward tasks may cause {Guanaco} to err if it does not decompose them into sequential logical steps—an established shortcoming \cite{weichain}. The following dialogue illustrates this issue:

\begin{quote}
    \textbf{User}: Please factorize 1833.

    \textbf{Guanaco}: \texttt{Sure, I'd be happy to help! 1833 is a prime number, which means it has only two factors: 1 and 1833. \\ \\ So the unique factorization of 1833 is: \\ \\ 2\^{}1 * 3\^{}2 * 17\^{}1 \\ \\ I hope this helps! Let me know if you have any other questions.}
\end{quote}

which is incorrect (the correct factorization is $3 \times 17 \times 43$), and in fact, doubly so.

\paragraph{Theory of Mind} {Guanaco} exhibits unexpectedly advanced Theory of Mind reasoning \cite{nematzadeh2018evaluating, sap2022neural}. For example, the model provides a thorough and accurate response to the following prompt:

\begin{quote}
    \textbf{User}: James and Abby are in the bedroom. Abby put the pen in the desk drawer. Abby leaves the bedroom. James moves the pen into the bag. Where does James think Abby will look for the pen?

    \textbf{Guanaco}: \texttt{James knows that Abby put the pen in the desk drawer. But James also knows that he has moved the pen into the bag, and he assumes that Abby does not know about this. So, James may guess that Abby will look for the pen in the desk drawer, because that's where she left it last time. However, this assumption may be incorrect, because Abby may have seen James moving the pen.}
\end{quote}

Nonetheless, these types of inferences are not consistently reliable, as the model sometimes gives explanations based on unjustified assumptions, such as:

\begin{quote}
    \textbf{User}: Evelyn entered the living room. Jackson entered the playroom. James entered the playroom. The beans are in the treasure chest. James exited the playroom. Jackson moved the beans to the pantry. Jackson exited the playroom. James entered the living room. Where will James look for the beans?

    \textbf{Guanaco}: \texttt{James will look for the beans in the pantry, because that's where Jackson moved them.}
\end{quote}

In this instance, {Guanaco} incorrectly assumes knowledge transfer that was never mentioned. Such shortcomings align with observations in recent research \cite{sap2022neural}, highlighting the need for further exploration.

\subsection{Considerations}
\label{ssec:considerations}

\paragraph{Evaluation}

Our findings show only moderate agreement among human annotators (Fleiss $\kappa=0.42$), with the level of agreement dropping further when contrasting outputs from two advanced systems. This underscores the insufficiency of existing evaluation benchmarks and human judgment protocols for assessing chatbot effectiveness. In head-to-head evaluations of ChatGPT and {Guanaco} 65B outputs on the Vicuna benchmark, we notice that individual preferences significantly influence outcomes, as co-authors of this study frequently differed in their chosen responses. Advancing evaluation frameworks will require drawing on fields like Human-Computer Interaction and Psychology, where strategies exist for reconciling subjective judgments.

Additionally, we discovered notable biases in automated evaluation methods. For example, GPT-4 consistently scores the model appearing first in its prompt more favorably, demonstrating a strong order effect. The relatively low sample-level concordance between GPT-4 and human raters (Fleiss $\kappa=0.25$) further suggests a mismatch between machine-generated and human-derived assessments. Moreover, as seen in Table~\ref{tbl:elo_large}, GPT-4 tends to rate its own completions considerably higher than humans do, with Elo scores of 1348 compared to 1176—corresponding to an increased likelihood of winning of roughly 20\% over its competitor.

Future investigations should focus on identifying and addressing potential biases inherent in automated evaluation frameworks, as well as exploring methods for mitigating such biases.

\paragraph{Data \& Training} The OASST1 dataset utilized to train the {Guanaco} models features multilingual data, and the OA benchmark similarly contains prompts in multiple languages. The extent to which this multilingual training influences model effectiveness on instructions issued in non-English languages, and whether it accounts for the more substantial performance difference observed between the Vicuna-13B model (trained exclusively on English data) and the {Guanaco} 33B and 65B models on the OA benchmark, remains an open question for future research.

Given the notable performance of the {Guanaco} models, we conducted an analysis for potential data contamination between the OASST1 dataset and prompts from the Vicuna benchmark. Utilizing fuzzy string matching alongside manual review of the most similar pairs, we did not detect any overlapping prompts between these two data sources.

Additionally, it is important to highlight that our models are trained solely with cross-entropy loss in a supervised learning framework and do not utilize reinforcement learning from human feedback (RLHF). This observation prompts further research into understanding the comparative benefits and drawbacks of training with only cross-entropy loss versus incorporating RLHF. We anticipate that \textsc{QLoRA} can facilitate such comprehensive analyses at scale without demanding excessive computational resources.

\section{Related Work}

\paragraph{Quantization of Large Language Models} Most prior research on LLM quantization has centered around inference-time quantization. Key methods aimed at preserving the performance of 16-bit models emphasize the treatment of outlier features, as in SmoothQuant~\citep{xiao2022smoothquant} and LLM.int8()~\citep{dettmers2022llm}, whereas others leverage advanced grouping strategies \citep{park2022nuqmm, yao2022zeroquant}. Lossy quantization research explores the compromises associated with standard rounding techniques~\citep{dettmers2022case,zeng2022glm,pope2022efficiently} and investigates methods to optimize rounding in order to achieve higher quantization precision~\citep{frantar2022gptq}. Apart from our present study, only SwitchBack layers~\citep{wortsman2023stable} has examined backpropagation through quantized weights at scales greater than 1B parameters.

\paragraph{Finetuning with Adapters} In our experiments, we employ Low-rank Adapters~\citep{hu2021lora} (LoRA), though the landscape of Parameter Efficient FineTuning (PEFT) techniques is broad. Alternatives include prompt tuning~\citep{qin2021learning,lester2021power,li2021prefix}, modifying the inputs to embedding layers~\citep{an2022input}, tuning hidden states (as in IA$^3$)~\citep{liu2022few}, inserting entire additional layers~\citep{houlsby2019parameter}, adjusting only the bias parameters~\citep{zaken2021bitfit}, masking weights based on Fisher information~\citep{sung2021training}, and various hybrid strategies~\citep{henderson2021compacter}. Our results indicate that LoRA adapters are capable of attaining the same performance as full 16-bit finetuning. The comparative merits of other PEFT methods remain to be explored in subsequent research.

\paragraph{Instruction Finetuning} Instruction finetuning aims to enable a pretrained LLM to more effectively follow prompts by adapting the model on input-output pairs spanning diverse data sources, teaching the LLM to produce appropriate outputs conditioned on prompts. Related strategies and benchmarks include MetaICL~\citep{min2021metaicl}, MetaTuning~\citep{zhong2021adapting}, InstructGPT~\citep{ouyang2022training}, FLAN~\citep{wei2021finetuned,chung2022scaling}, PromptSource~\citep{bach2022promptsource}, Super-NaturalInstructions~\citep{wang2022super,sanh2021multitask}, Self-instruct~\citep{wang2022self}, UnnaturalInstructions~\citep{honovich2022unnatural}, OPT-IML~\citep{iyer2022opt}, UnifiedSKG~\citep{xie2022unifiedskg}, OIG/Chip2~\citep{laion2023}, Alpaca~\citep{alpaca}, Vicuna~\citep{vicuna2023}, Koala~\citep{koala_blogpost_2023}, and Self-instruct-GPT-4~\citep{peng2023instruction}.

\paragraph{Chatbots} Instruction-tuned models are frequently presented as conversational chatbots, with training often incorporating Reinforcement Learning from Human Feedback (RLHF)~\citep{christiano2017deep} or alternative procedures such as utilizing AI feedback for model improvement (RLAIF)~\citep{bai2022constitutional}. Notable models and resources in this space include Anthropic-HH~\citep{askell2021general,bai2022training}, Open Assistant~\citep{kopf2023openassistant}, LaMDA~\citep{thoppilan2022lamda}, and Sparrow~\citep{glaese2022improving}. Our approach does not leverage reinforcement learning. Instead, Guanaco—our leading model—is finetuned using multi-turn conversational data from Open Assistant, a dataset purpose-built for RLHF-style training~\citep{kopf2023openassistant}. For chatbot evaluation, some recent approaches have adopted GPT-4-based assessments in lieu of expensive human annotation~\citep{vicuna2023,peng2023instruction}. We build upon and improve these evaluation protocols, placing an emphasis on more dependable assessment methodology.

\section{Limitations and Discussion}

Our findings demonstrate that \textsc{QLoRA} can, with a 4-bit base and Low-rank Adapters (LoRA), achieve performance comparable to that of full 16-bit finetuning. Nevertheless, it remains to be established whether this equivalence persists for models of 33B and 65B parameters; we defer the empirical analysis at these larger scales due to prohibitive computational requirements.

An additional constraint of this work concerns the scope of our evaluation for instruction finetuned models. Our reported results are based on the MMLU benchmark, the Vicuna benchmark, and the OA benchmark; however, evaluations have not been extended to alternative benchmarks like BigBench, RAFT, and HELM, raising questions about the generality of our conclusions. That said, we present a comprehensive study on MMLU and introduce novel evaluation methodologies for chatbots.

Based on the data discussed, benchmark outcomes seem heavily influenced by the degree of alignment between the finetuning data and the specific benchmark tasks. For instance, FLAN v2 shares similarities with the MMLU dataset but differs from chatbot-related evaluations, whereas the Chip2 dataset shows the opposite trend, and the scores for both models reflect this distinction on the MMLU and Vicuna benchmarks. This observation underscores not only the need for improved benchmarking and evaluation methods, but also the importance of carefully selecting the tasks under scrutiny. It prompts us to reflect on our priorities: should the goal be to optimize for knowledge and reasoning skills assessed in high school or college-like settings, or for conversational competence as seen in chatbots? Or is there yet another target capability? Since there is a natural tendency to reuse established benchmarks rather than to design novel ones, these benchmarks can influence research trajectories. It is thus imperative for the research community to ensure that benchmark metrics genuinely capture the capabilities of greatest interest.

While this work presents a comprehensive assessment of general chatbot performance, one recognized shortcoming is the relatively limited scope of the responsible AI analysis performed for {Guanaco}. Specifically, we analyzed the propensity of {Guanaco}-65B to generate socially biased outputs in comparison to other models (see Table~\ref{tbl:crows}). The results show that the average bias score for {Guanaco}-65B is considerably lower than that observed in raw, pretrained counterparts, indicating that finetuning on the OASST1 dataset mitigates some biases present in the LLaMA base model. Although this outcome is promising, the extent to which {Guanaco} addresses other dimensions of bias remains uncertain, and a more exhaustive bias evaluation is left for subsequent studies.

Another limitation involves the lack of analysis on alternative quantization levels, such as 3-bit base models, as well as various adapter methods. While this work employed LoRA, which is widely recognized for its robustness, there exists a broader array of Parameter Efficient FineTuning (PEFT) strategies with demonstrated success. It is not currently clear if such alternatives can be effectively scaled to models of significant size. Given that finetuning following quantization often recovers a substantial portion of information lost during the process, this suggests the possibility of even more aggressive quantization. As an illustrative example, employing 3-bit GPTQ quantization alongside LoRA on the basemodel may allow the finetuned model to match the performance achieved by full 16-bit finetuning.

\section{Broader Impacts}

The \textsc{QLoRA} finetuning technique represents the first approach capable of finetuning 33B parameter models using a single consumer-grade GPU, and 65B parameter models on a single professional-grade GPU, all without a loss of performance relative to standard full-parameter finetuning. Our findings demonstrate that the best 33B parameter model, when trained with the Open Assistant dataset, can perform comparably to ChatGPT on the Vicuna benchmark. As instruction finetuning is pivotal for adapting generic pretrained LLMs into ChatGPT-style conversational agents, we anticipate that this technique will facilitate much broader access to finetuning, particularly benefiting researchers and practitioners with limited computational resources, thus enhancing the accessibility of leading-edge NLP tools. The introduction of \textsc{QLoRA} acts as a leveling mechanism, helping to reduce the disparity between resource-rich organizations and smaller teams operating with consumer hardware.

An additional avenue for significant impact lies in the deployment of models to mobile devices. We contend that our \textsc{QLoRA} technique represents a breakthrough in enabling the finetuning of LLMs directly on mobile phones and in other environments with limited computational resources. Although it has previously been demonstrated that 7B parameter models can be executed on smartphones, \textsc{QLoRA} is the first approach to facilitate the finetuning process for these models in such settings. Our estimates suggest that, using an iPhone 12 Plus, \textsc{QLoRA} can process up to 3 million tokens overnight during charging periods. Despite the fact that finetuned 7B models do not match the performance of ChatGPT, we argue that their capabilities suffice to power innovative applications that were previously unfeasible due to concerns related to privacy or the effectiveness of LLMs. By supporting user-managed data and models, \textsc{QLoRA} enables more privacy-conscious LLM use cases, while also reducing deployment barriers for such systems.

Nevertheless, it is important to acknowledge that finetuning represents a dual-use capability that may be misused with harmful consequences. The risks associated with the widespread proliferation of LLMs are well documented \citep{bommasani2021opportunities,bender2021dangers}; however, we maintain that democratizing access to rapidly evolving technologies fosters broader and more independent oversight than restricting LLM capabilities to a few large organizations who may not provide the requisite transparency through open models or accessible source code.

In summary, we anticipate that \textsc{QLoRA} will substantially increase the availability and accessibility of high-quality LLM finetuning, thereby producing a largely beneficial effect.